% ==== Responsável: Mariana Nunes Pignaton ====
\chapter{Memória secundária e partições}\label{cap:particoes}\label{cap:memoria-sec}

A memória secundária é responsável pelo armazenamento permanente de dados e programas.
Diferentemente da memória principal (RAM), ela mantém as informações mesmo com o computador
desligado. Nela ficam o sistema operacional, os aplicativos e os arquivos dos usuários. Em
troca, o acesso é bem mais lento que o da RAM, por isso o sistema operacional precisa
organizar e escalonar o uso dos discos \cite{silberschatz2015}.

\section{Tipos de dispositivos de armazenamento}\label{sec:tipos-armazenamento}

\begin{itemize}
    \item \textbf{Disco rígido (HD):} usa discos magnéticos que giram, com cabeçotes de leitura
    e escrita que se movem até a trilha desejada. É mais lento que o SSD, mas continua sendo
    usado pela grande capacidade e pelo custo menor por \textit{gigabyte}.
    \item \textbf{Unidade de estado sólido (SSD):} usa memória \textit{flash}, sem partes
    mecânicas. Por não ter tempo de busca nem de rotação, é muito mais rápido que o HD, o que
    acelera o \textit{boot} e a abertura de programas.
    \item \textbf{Dispositivos removíveis:} \textit{pendrives}, cartões de memória e HDs
    externos, usados para transportar arquivos e fazer cópias de segurança.
\end{itemize}

\section{Particionamento de discos}\label{sec:particionamento}

O particionamento divide um dispositivo de armazenamento em partes independentes, chamadas
partições. Cada partição pode ter o seu próprio sistema de arquivos e ser usada para guardar
dados ou para instalar um sistema operacional. Isso permite, por exemplo, separar os arquivos
pessoais (\texttt{/home}) dos arquivos do sistema, ou instalar dois sistemas operacionais no
mesmo disco \cite{tanenbaum2016}.

No esquema MBR existem três tipos de partição:

\begin{itemize}
    \item \textbf{primária:} pode guardar um sistema operacional ou arquivos do usuário;
    \item \textbf{estendida:} funciona como um contêiner que permite criar partições lógicas;
    \item \textbf{lógica:} é criada dentro de uma partição estendida.
\end{itemize}

\subsection{Esquemas MBR e GPT}\label{subsec:mbr-gpt}

MBR e GPT são os dois padrões usados para guardar a tabela de partições de um disco. A
\autoref{tab:mbr-gpt} compara os dois.

\begin{table}[H]
    \centering
    \caption{Comparação entre os esquemas de partição MBR e GPT}
    \label{tab:mbr-gpt}
    \begin{tabularx}{\linewidth}{lXX}
        \toprule
        \textbf{Característica} & \textbf{MBR} & \textbf{GPT} \\
        \midrule
        Onde fica a tabela & Primeiro setor do disco (512 bytes) & Início do disco, com cópia de segurança no final \\
        Partições & Até 4 primárias (ou 3 primárias e 1 estendida) & 128 por padrão \\
        Tamanho máximo do disco & Cerca de 2\,TiB & Acima de 2\,TiB \\
        \textit{Firmware} típico & BIOS & UEFI \\
        \bottomrule
    \end{tabularx}
    \legend{Fonte: os autores (2026), com base em \citeonline{tanenbaum2016}.}
\end{table}

O MBR é o padrão mais antigo e tem limites de tamanho de disco e de número de partições. O GPT
é o padrão atual dos computadores com UEFI (ver \autoref{tab:bios-uefi}, no
\autoref{cap:boot}).

\section{Gerenciamento de partições no Linux}\label{sec:particoes-linux}

O Linux tem comandos para ver e administrar os dispositivos de armazenamento e as suas
partições, reunidos no \autoref{lst:particoes}. O \texttt{lsblk}\footnote{Manual do comando
\texttt{lsblk}. Disponível em: \url{https://man7.org/linux/man-pages/man8/lsblk.8.html}.
Acesso em: 20 nov. 2026.} mostra os discos e as partições em forma de árvore (por exemplo,
\texttt{sda}, \texttt{sda1}, \texttt{sda2}). O \texttt{fdisk -l} mostra a tabela de partições
de cada disco. O \texttt{df -h} mostra o espaço usado e livre em cada sistema de arquivos
montado, e o \texttt{du -sh} mostra quanto espaço um diretório ocupa.

\begin{lstlisting}[caption={Comandos para consultar discos e partições no Linux},label={lst:particoes}]
lsblk                  # discos e particoes em arvore
sudo fdisk -l          # tabela de particoes de cada disco
df -h                  # espaco usado e livre em cada sistema de arquivos
du -sh ~/Documentos    # espaco ocupado por um diretorio
\end{lstlisting}

O comando \texttt{df -h} também foi usado na opção 1 do menu da Atividade 7.11
(\autoref{fig:a711a}), em que aparecem a partição \texttt{/dev/sda2}, montada na raiz
(\texttt{/}), e sistemas de arquivos temporários na memória (\texttt{tmpfs}).
